\section{Attention Head 的典型模式}

\subsection{常见的 Head 模式}

通过构造测试序列并可视化不同 Head 的 Attention 矩阵，可以识别出以下典型模式：

\begin{importantbox}{五种典型 Attention 模式}
\begin{enumerate}
  \item \textbf{Local Head}：主要关注相邻 token，呈对角线形式
  \item \textbf{Strided Head}：以固定步长关注 token，捕获周期性模式
  \item \textbf{Global Head}：几乎均匀关注所有 token
  \item \textbf{Position Head}：关注特定位置（如句首、段落开头）
  \item \textbf{Sink Head}：将大量注意力集中在序列的第一个 token 上
\end{enumerate}
\end{importantbox}

\subsection{不同层的功能分工}

一般规律：
\begin{itemize}
  \item \textbf{浅层}（底层）：更多 Local Head，处理语法、词组搭配
  \item \textbf{中层}：混合模式，开始建立语义关联
  \item \textbf{深层}（顶层）：更多 Global 和 Position Head，处理高层语义
\end{itemize}

\subsection{本章小结}

不同 Head 自然分化出不同功能。浅层关注局部语法，深层关注全局语义，形成了一个从底向上的抽象层次结构。

